%!TEX program = xelatex
\documentclass{article}
\usepackage{graphicx} % Required for inserting images

% ─── Packages ───────────────────────────────────────────────
\usepackage{fontspec}
\usepackage{polyglossia}
\usepackage{xcolor}
\usepackage{colortbl}
\usepackage{booktabs}
\usepackage{array}
\usepackage{geometry}
\usepackage{tabularx}
\usepackage{amssymb}

% ─── Page geometry ──────────────────────────────────────────
\geometry{
  a5paper,
  margin=1cm
}

% ─── Language & Font setup ──────────────────────────────────
\setmainlanguage{bengali}
\setotherlanguage{arabic}

\newfontfamily\bengalifont[
  Script=Bengali,
  BoldFont=SolaimanLipi
]{SolaimanLipi}

\newfontfamily\arabicfont[
  Script=Arabic,
  Scale=1.6,
  WordSpace=1.5
]{Scheherazade New}

% ─── Colors ─────────────────────────────────────────────────
\definecolor{headergreen}{HTML}{4CAF50}
\definecolor{lightgreen}{HTML}{E8F5E9}
\definecolor{suffixred}{HTML}{CC0000}
\definecolor{titlegreen}{HTML}{2E7D32}

% ─── Helper commands ────────────────────────────────────────
\newcommand{\myroot}[1]{{\color{black}#1}}
\newcommand{\sfx}[1]{{\color{suffixred}#1}}

\newcolumntype{C}{>{\centering\arraybackslash}X}

% ─── Arabic words ───────────────────────────────────────────
% Each word is wrapped in \textarabic{} so the whole word is ONE
% right-to-left unit. Without this, the coloured root and suffix
% get reordered left-to-right and the red lands on the wrong side.
%
% SPECIAL CASE — when the marker (علامة) is a bare diacritic, not an
% added letter (e.g. the final fatha in فَعَلَ for واحد مذكر غائب):
% XeLaTeX cannot colour a diacritic differently from its base letter,
% because changing colour mid-cluster detaches the mark. Instead we
% OVERLAY: the fully-red word is drawn first (bottom layer, and it sets
% the box width), then a black copy LACKING the final mark is \llap-ed
% on top. The black letters hide the red ones; only the red diacritic
% shows through. Kept inside \mbox so it centres correctly in a cell.
% Usage: \redmark{full red word}{same word WITHOUT the final mark}
\newcommand{\redmark}[2]{\mbox{\textarabic{\sfx{#1}}\llap{\textarabic{\myroot{#2}}}}}

% غائب
\newcommand{\fala}     {\redmark{فَعَلَ}{فَعَل}}
\newcommand{\falaa}    {\textarabic{\myroot{فَعَلَ}\sfx{ا}}}
\newcommand{\faluu}    {\textarabic{\myroot{فَعَلُ}\sfx{وا}}}
\newcommand{\falat}    {\textarabic{\myroot{فَعَلَ}\sfx{تْ}}}
\newcommand{\falataa}  {\textarabic{\myroot{فَعَلَ}\sfx{تَا}}}
\newcommand{\falna}    {\textarabic{\myroot{فَعَلْ}\sfx{نَ}}}

% حاضر
\newcommand{\falta}    {\textarabic{\myroot{فَعَلْ}\sfx{تَ}}}
\newcommand{\faltumaa} {\textarabic{\myroot{فَعَلْ}\sfx{تُمَا}}}
\newcommand{\faltum}   {\textarabic{\myroot{فَعَلْ}\sfx{تُمْ}}}
\newcommand{\falti}    {\textarabic{\myroot{فَعَلْ}\sfx{تِ}}}
\newcommand{\faltumaaf}{\textarabic{\myroot{فَعَلْ}\sfx{تُمَا}}}
\newcommand{\faltunna} {\textarabic{\myroot{فَعَلْ}\sfx{تُنَّ}}}

% متكلم
\newcommand{\faltu}    {\textarabic{\myroot{فَعَلْ}\sfx{تُ}}}
\newcommand{\falnaa}   {\textarabic{\myroot{فَعَلْ}\sfx{نَا}}}

% ─── المضارع (present/future) ────────────────────────────────
% In المضارع the markers are PREFIXES (حروف المضارعة: يَ تَ أَ نَ) as
% well as suffixes. So both ends are red, the root فعل stays black.
%
% \mud{prefix}{black stem}{suffix}: a form whose marker suffix is an
% ADDED LETTER (انِ، ونَ، ينَ، نَ). Straightforward coloured runs.
\newcommand{\mud}[3]{\textarabic{\sfx{#1}\myroot{#2}\sfx{#3}}}
%
% \mudraf{prefix}: a singular form ending in a bare damma (raf'), e.g.
% يَفْعَلُ. The final damma is the marker but a bare diacritic, so —as
% with the ماضي fatha— we OVERLAY. Crucially BOTH layers split the
% prefix the same way, so they shape (and measure) identically and
% align. Bottom = red prefix + red stem+damma; top = red prefix +
% black stem (no damma). Black hides red stem; red damma shows.
\newcommand{\mudraf}[1]{%
  \mbox{\textarabic{\sfx{#1}\sfx{فْعَلُ}}\llap{\textarabic{\sfx{#1}\myroot{فْعَل}}}}%
}

% ─── Document info ──────────────────────────────────────────
\title{সরফ}
\author{মাজহার আহমেদ}
\date{জুন ২০২৬}

\begin{document}

\maketitle

\section{ফেল এর প্রকারভেদ}

\begin{itemize}
  \item মাজি
  \item মুজারে
  \item আমর
  \item নাহি
  \item ইসমে মুসতাক
\end{itemize}

\newpage

% ─── Fil Madi Table ─────────────────────────────────────────
\begin{center}
  {\arabicfont \color{titlegreen} \textbf{\textarabic{تصريف فِعل الماضي}}}
\end{center}

\vspace{0.3cm}

\begin{center}
  {\arabicfont
    \textarabic{%
      \textbf{\color{titlegreen}علامة:}
      \enspace
      الجذر \textbf{\myroot{فَعَلَ}}
      \enspace|\enspace
      الضمير (علامة) \textbf{\color{suffixred}بالأحمر}%
    }
  }
\end{center}

\vspace{0.4cm}

\setlength{\arrayrulewidth}{0.5pt}
\setlength{\tabcolsep}{3pt}
\renewcommand{\arraystretch}{2.2}

\begin{tabularx}{\textwidth}{|C|C|C|}

  \hline
  \rowcolor{headergreen}
  {\arabicfont \color{white} \textbf{مُتَكَلِّم}} &
  {\arabicfont \color{white} \textbf{حَاضِر}}     &
  {\arabicfont \color{white} \textbf{غَائِب}}     \\
  \hline

  \rowcolor{lightgreen}
  {\arabicfont \faltu}   &
  {\arabicfont \falta}   &
  {\arabicfont \fala}    \\
  \hline

  {\arabicfont \falnaa}     &
  {\arabicfont \faltumaa}   &
  {\arabicfont \falaa}      \\
  \hline

  \rowcolor{lightgreen}
  &
  {\arabicfont \faltum}  &
  {\arabicfont \faluu}   \\
  \hline

  &
  {\arabicfont \falti}   &
  {\arabicfont \falat}   \\
  \hline

  \rowcolor{lightgreen}
  &
  {\arabicfont \faltumaaf} &
  {\arabicfont \falataa}   \\
  \hline

  &
  {\arabicfont \faltunna} &
  {\arabicfont \falna}    \\
  \hline

\end{tabularx}

\vspace{0.4cm}

\begin{center}
  {\arabicfont \color{gray} \small
    \textarabic{الجذر الأصلي بالأسود \textbullet\ العلامة (الضمير) بالأحمر}
  }
\end{center}

% ─── NEW PAGE: Fil Mudare Table ─────────────────────────────
\newpage

\begin{center}
  {\arabicfont \color{titlegreen} \textbf{\textarabic{تصريف فِعل المضارع}}}
\end{center}

\vspace{0.3cm}

\begin{center}
  {\arabicfont
    \textarabic{%
      \textbf{\color{titlegreen}علامة:}
      \enspace
      الجذر \textbf{\myroot{فعل}}
      \enspace|\enspace
      حرف المضارعة والعلامة \textbf{\color{suffixred}بالأحمر}%
    }
  }
\end{center}

\vspace{0.4cm}

\setlength{\arrayrulewidth}{0.5pt}
\setlength{\tabcolsep}{3pt}
\renewcommand{\arraystretch}{2.2}

\begin{tabularx}{\textwidth}{|C|C|C|}

  \hline
  \rowcolor{headergreen}
  {\arabicfont \color{white} \textbf{مُتَكَلِّم}} &
  {\arabicfont \color{white} \textbf{حَاضِر}}     &
  {\arabicfont \color{white} \textbf{غَائِب}}     \\
  \hline

  % Row 1 — واحد مذكر
  \rowcolor{lightgreen}
  {\arabicfont \mudraf{أَ}}            &
  {\arabicfont \mudraf{تَ}}            &
  {\arabicfont \mudraf{يَ}}            \\
  \hline

  % Row 2 — تثنية مذكر / جمع متكلم
  {\arabicfont \mudraf{نَ}}            &
  {\arabicfont \mud{تَ}{فْعَلَ}{انِ}}  &
  {\arabicfont \mud{يَ}{فْعَلَ}{انِ}}  \\
  \hline

  % Row 3 — جمع مذكر
  \rowcolor{lightgreen}
  &
  {\arabicfont \mud{تَ}{فْعَلُ}{ونَ}}  &
  {\arabicfont \mud{يَ}{فْعَلُ}{ونَ}}  \\
  \hline

  % Row 4 — واحد مؤنث
  &
  {\arabicfont \mud{تَ}{فْعَلِ}{ينَ}}  &
  {\arabicfont \mudraf{تَ}}            \\
  \hline

  % Row 5 — تثنية مؤنث
  \rowcolor{lightgreen}
  &
  {\arabicfont \mud{تَ}{فْعَلَ}{انِ}}  &
  {\arabicfont \mud{تَ}{فْعَلَ}{انِ}}  \\
  \hline

  % Row 6 — جمع مؤنث
  &
  {\arabicfont \mud{تَ}{فْعَلْ}{نَ}}   &
  {\arabicfont \mud{يَ}{فْعَلْ}{نَ}}   \\
  \hline

\end{tabularx}

\vspace{0.4cm}

\begin{center}
  {\arabicfont \color{gray} \small
    \textarabic{حرف المضارعة (يَ تَ أَ نَ) والعلامة بالأحمر \textbullet\ الجذر بالأسود}
  }
\end{center}

\end{document}
